% !TeX program = lualatex
% halo:begin
% {
%   "version": 1,
%   "publish": false,
%   "course": "data-and-optimization",
%   "section": "1.1",
%   "title": "数据与优化（通识笔记）1.1：大数据与最优化",
%   "categories": ["study-notes"],
%   "tags": ["mathematics"],
%   "excerpt": "大数据基础与思维、机器学习和深度学习、优化建模、回归与分类、稀疏与低秩结构，以及优化算法与收敛。"
% }
% halo:end
% 来源：通识课-1大数据案例介绍.pdf（1--60页）；通识课-2大数据思维.pdf（1--33页）；
% 01-opt-dzw.pdf（1--45页）；05-lect1-model.pdf（1--31页，支持向量机结束）。
% 课程范围由用户于2026-10-28确认；教材不纳入复习范围。
% 来源定位与写作阶段记录见Beads SHU-notes-0w4，修订版本2。
% 2026-11-03：按读者反馈补充解释、教学例子与开卷查阅入口。
% 信息论公式、统一符号及简短推导用于解释课件概念，不另扩展教材章节。
\RequirePackage{iftex}
\RequireLuaTeX
\documentclass[UTF8,fontset=fandol,a4paper,12pt]{ctexart}
\usepackage[margin=2.5cm]{geometry}
\usepackage{amsmath,amssymb,amsthm}
\usepackage{enumitem,fancyhdr,booktabs,longtable,array}
\usepackage[hidelinks]{hyperref}
\newcommand{\R}{\mathbb R}
\newcommand{\E}{\mathbb E}
\newcommand{\rank}{\operatorname{rank}}
\newcommand{\vx}{\boldsymbol x}
\newcommand{\va}{\boldsymbol a}
\newcommand{\vb}{\boldsymbol b}
\newcommand{\vw}{\boldsymbol w}
\newcommand{\norm}[1]{\left\lVert#1\right\rVert}
\theoremstyle{definition}
\newtheorem{definition}{定义}[section]
\newtheorem{proposition}[definition]{命题}
\newtheorem{example}[definition]{例}
\renewcommand{\proofname}{证明}
\setlist[enumerate]{leftmargin=2em,itemsep=0.3em,topsep=0.3em}
\setlist[itemize]{leftmargin=2em,itemsep=0.2em,topsep=0.3em}
\setlength{\headheight}{16pt}
\pagestyle{fancy}\fancyhf{}
\fancyhead[L]{\small 数据与优化（通识笔记）}
\fancyhead[R]{\small 大数据与最优化}
\fancyfoot[C]{\thepage}
\renewcommand{\headrulewidth}{0.3pt}
\raggedbottom
\title{大数据与最优化}
\author{Yuchen Zhang}
\date{修订日期：2026年11月3日}
\makeatletter
\renewcommand{\maketitle}{\begin{center}
{\LARGE\bfseries\@title\par}\vspace{0.6em}
\@author\par\vspace{0.2em}{\small\@date\par}
\end{center}}
\makeatother
\begin{document}
\maketitle
评分记录可以用于推荐，搜索记录可以辅助预测，交通记录可以支持路线选择。数据的作用取决于所研究的问题，也取决于怎样从记录中提取信息、建立模型并检验结果。
\section*{关键词查阅}
说明一种方法，可以沿着“解决什么问题、已知什么与要求什么、怎样选择结果、结果有什么作用与限制”展开。公式中的每一项都应能对应到其中一个具体要求。
\begin{center}
\begin{tabular}{p{0.27\textwidth}p{0.51\textwidth}r}
\toprule
关键词 & 主要解释 & 页码\\
\midrule
数据、4V、数据价值 & 数据量大为什么还需要有意义、可靠的记录 & \pageref{sec:data}\\
不确定性、熵、相关因果 & 信息怎样帮助判断，为什么相关还不能证明因果 & \pageref{sec:thinking}\\
推荐、预测、隐私 & 怎样用同一流程解释不同案例 & \pageref{sec:cases}\\
学习、表示、泛化、网络 & 训练已经很好，为什么新数据仍可能预测不好 & \pageref{sec:ml}\\
变量、目标、约束、松弛 & 投资要求怎样成为公式，改变写法是否改变问题 & \pageref{sec:model}\\
回归、似然、最小二乘 & 系数怎样从观测中估计 & \pageref{sec:regression}\\
正则、岭、LASSO & 为什么限制系数；缩小与筛选的差别 & \pageref{sub:regularization}\\
稀疏、$\ell_0/\ell_1/\ell_2$ & 多个可行解之间怎样选 & \pageref{sec:sparse}\\
低秩、核范数、矩阵补全 & 怎样利用共同偏好预测缺失评分 & \pageref{sec:lowrank}\\
逻辑回归、SVM、软间隔 & 概率分类与最大间隔分类 & \pageref{sec:classification}\\
梯度、SGD、牛顿、动量 & 更新方向、步长与计算代价 & \pageref{sec:algorithms}\\
最优、收敛、速度、停止 & 怎样评价算法算出的结果 & \pageref{sec:convergence}\\
\bottomrule
\end{tabular}
\end{center}
\clearpage

\section{数据、信息与大数据}\label{sec:data}
\subsection{从记录到知识}
\begin{definition}
\textbf{数据}是对事物的状态、属性或变化的记录，可以表现为数值、文字、图像、声音和行为记录等。\textbf{信息}是数据在具体问题中表达的意义。\textbf{知识}是对信息进一步组织、解释后形成的系统认识。
\end{definition}

例如，天体在各个时刻的位置是数据，由这些记录得到的运动轨迹是信息，对运动规律的概括则形成知识。由观测建立模型，再用模型预测新的观测，构成了从数据取得认识的一条基本途径。

同样多的记录，所能提供的信息未必相同。重复一条已知记录，通常不会增加同等数量的新信息；与研究目标无关的数据，也未必能改善判断。因此，数据的价值既取决于记录的内容和质量，也取决于所要解决的问题。

\subsection{大数据的特点与来源}
大数据的困难不仅在于数量多，还在于数据产生快、类型多，并且有用信息需要从大量记录中提取。其基本特点常概括为\textbf{4V}。
\begin{center}
\begin{tabular}{p{0.23\textwidth}p{0.66\textwidth}}
\toprule
特点 & 含义及处理要求\\
\midrule
规模巨大（Volume） & 数据量超出常规处理方式的有效承受范围，需要考虑存储、计算与扩展能力。\\
高速产生（Velocity） & 数据持续生成、更新，分析结果常有时效要求，迟到的结果可能失去作用。\\
形式多样（Variety） & 数值表格、文本、图像、音频、视频等并存，需要不同的表示和处理方法。\\
潜藏价值（Value） & 大量记录中蕴含有用规律，但价值密度可能较低，需要围绕任务进行分析。\\
\bottomrule
\end{tabular}
\end{center}

数据库中的规则表格属于结构化数据，文章、影像和语音通常属于非结构化数据。传感器、移动终端、社交平台和交易系统不断产生这些记录，移动互联网又使不同记录能够相互关联。

\textbf{大数据与单纯的大规模数据有所区别。}除了体量，还要考察数据是否具有多维关联，是否覆盖研究所需的情形，以及能否及时更新。例如，同一条交通记录同时包含位置、时间、速度和行驶方向，就能支持比单独统计车辆总数更丰富的判断。这里的完备性指与任务有关的覆盖程度，并不意味着掌握了所有事实。

\subsection{数据分析的完整过程}
从数据到结论，需要经过问题定义、数据处理、建模和验证。
\begin{enumerate}
\item \textbf{确定问题。}明确要描述什么、预测什么或选择什么，并规定评价标准。
\item \textbf{收集与处理数据。}检查来源、缺失、重复、异常和测量误差，统一记录的时间、单位与口径。
\item \textbf{形成表示。}把原始记录转化为模型能够利用的样本、特征或矩阵。
\item \textbf{选择模型并估计参数。}从较简单的模型起步，用数据判断是否有必要增加复杂性。
\item \textbf{验证与比较。}在相同条件下比较候选方法，检验新数据上的误差、计算代价与适用范围。
\item \textbf{应用与反馈。}把结果用于预测或决策，并在环境变化后重新检查模型。
\end{enumerate}
数据收集的\textbf{量}与\textbf{质}必须同时考虑。错误记录和缺乏代表性的样本，即使数量很多，也可能使后续分析偏离实际目标。

在这一过程中，数据库技术负责组织、存储和查询记录，统计学研究数据中的变异、估计与推断，数据挖掘侧重发现有用模式，机器学习侧重从数据学习预测或表示。这些领域相互交叉；数据挖掘的英文为 data mining，深度学习为 deep learning，两者含义不同。

\section{大数据思维与不确定性}\label{sec:thinking}
\subsection{确定性描述与概率描述}
确定性模型用规律描述变量之间的关系；当影响因素复杂或观测不完整时，概率模型用分布表达可能结果。

机械思维强调：规律是确定的，可以用较简单的公式或语言描述，并能在相应条件下用于预测。其优势是结构清楚、解释明确；但当变量很多、关系复杂、观测不充分时，直接写出完整的确定性规律可能十分困难。

大数据思维承认判断中的不确定性，利用大量观测、多维关联和统计规律改善预测。\textbf{数据驱动}是从样本中学习规律或参数；\textbf{机理驱动}是从对象的内在机制建立关系。二者可以结合，例如用物理规律限制模型，再用数据估计未知参数。数据驱动不要求放弃解释，概率描述也不等于现象没有规律。

\subsection{信息熵与信息的作用}
判断一条道路是否拥堵，事先可能有多个结果；取得车辆速度等记录后，判断才有依据。信息的作用可以从这种不确定性的变化来理解。设离散随机变量 $X$ 的可能取值为 $x_1,\ldots,x_s$，相应概率为 $p_1,\ldots,p_s$。
\begin{definition}[信息熵]
\[
 \boxed{H(X)=-\sum_{i=1}^s p_i\log_2 p_i,}
\]
其中约定 $0\log_2 0=0$。以 $2$ 为底时，单位为比特。
\end{definition}
公式可以分两步读：$-\log_2p_i$ 表示结果 $x_i$ 发生时带来的信息量，越少见的结果带来的意外越大；再乘它出现的概率 $p_i$ 并求和，就得到事先平均有多少不确定性。例如抛一枚公平硬币，两面各有 $1/2$ 的概率，每个结果的信息量都是 $1$，加权平均也为 $1$ 比特。若事先知道一定正面朝上，则实际结果不会带来新的意外。
若结果已经确定，则熵为 $0$；若 $s$ 个结果等可能，则熵为 $\log_2 s$，也是固定取值个数下的最大熵。熵越大，事先判断结果的不确定性越大。

取得额外观测 $Y$ 后，对每一种可能的观测结果，分别计算关于 $X$ 剩下的不确定性，再按这些观测出现的概率加权平均，得到条件熵 $H(X\mid Y)$。它不是指每次观测都减少相同的量，而是考察平均效果。互信息
\[
 I(X;Y)=H(X)-H(X\mid Y)\ge0
\]
表示观测 $Y$ 平均减少了多少关于 $X$ 的不确定性。若 $Y$ 与 $X$ 独立，则互信息为 $0$；若新增记录只是已有信息的重复，通常也不能提供同等数量的新信息。例如，判断道路是否拥堵时，增加车辆位置和速度记录可能有用，反复保存同一时刻的同一位置则主要增加冗余。

这里要比较的是\textbf{观测与待判断对象}的关系。新增数据与旧数据独立，并不表示它对目标无用；两批独立抽取的样本仍可能共同改善总体均值的估计。数据有没有价值，要看它能否为目标提供新的有效信息。

\subsection{编码、传输与概率模型}
信息熵不仅解释不确定性，还给出了无损编码的基本界限。对于独立同分布的离散信源，采用可唯一译码的编码时，平均每个符号的码长不能低于信源熵；通过对足够长的符号块联合编码，可以使平均每个符号的码长任意接近这一界限。这是信源编码定理所说明的压缩极限。
这里“信源”是产生符号的来源，“独立同分布”指各次产生不相互影响且遵循同一分布，“可唯一译码”指收到编码后可以无歧义地还原符号。定理给出的并非某个现成压缩程序，而是在这些前提下压缩能达到的理论界限。

信道编码定理讨论有噪声的传输。在相应的信道模型下，低于信道容量的传输速率可以通过适当编码实现任意低的错误率；超过容量则不能保证可靠传输。\textbf{前者讨论怎样压缩，后者讨论怎样可靠传送}，两者所回答的问题不同。

当只掌握某些统计条件时，可以在满足这些条件的概率分布中选择熵最大的一个，这称为\textbf{最大熵原理}。它保留已知信息，同时避免对未知部分施加额外的无依据限制；并不表示可以忽略已有知识。
例如，只知道结果有两个可能取值，却没有偏向哪一个的其他条件时，等概率分布具有最大熵；若已经知道某个结果的实际概率，就必须满足这一条件，不能仍机械地取等概率。

比较实际分布与模型分布时，可以问：实际常出现的结果，模型有没有给它足够高的概率？交叉熵用实际出现的概率作权重，以模型概率的负对数作为代价。记实际分布为 $P=(p_i)$、模型分布为 $Q=(q_i)$，交叉熵与相对熵分别为
\[
 H(P,Q)=-\sum_i p_i\log_2 q_i,\qquad
 D_{\mathrm{KL}}(P\Vert Q)=\sum_i p_i\log_2\frac{p_i}{q_i}.
\]
在 $p_i>0$ 的位置要求 $q_i>0$，否则相应交叉熵为无穷大。两者满足
\[
 \boxed{H(P,Q)=H(P)+D_{\mathrm{KL}}(P\Vert Q).}
\]
因此，$P=Q$ 时相对熵为 $0$，而交叉熵等于 $H(P)$，一般并不为 $0$。例如两者都是 $(1/2,1/2)$ 时，交叉熵为 $1$ 比特，模型已经正确表达分布，但结果本身仍有不确定性。交叉熵可以衡量用模型分布解释数据所需的平均编码代价，后文分类模型的对数损失与此相联系。
相对熵则表示这项代价超出实际分布自身熵的部分，用来衡量模型分布与实际分布的差异。

\subsection{相关关系与因果关系}
相关关系描述共同变化，因果关系则涉及一个因素的改变怎样影响另一个因素。例如，某类搜索量与流感活动可能同时上升，因此搜索量能够成为预测线索；但人为增加搜索次数，并不意味着会增加患病人数。

两个变量相关，可能是一个影响另一个，也可能是反向作用，或二者受共同因素影响。因此，\textbf{相关关系可以支持预测，却不能单独证明因果关系}。要从预测转向干预，还需要相应的研究设计和证据。多来源记录之间的相互核验可以提高信息可信度，但“交叉核验”与后文用于评价模型的“交叉验证”是不同概念。

\subsection{数据量、代表性与误差}
增加样本量可以减小某些随机误差，却不能自动消除采样偏差。设随机变量 $X$ 有有限期望 $\mu$ 和方差 $\sigma^2$，则切比雪夫不等式给出
\[
 \Pr\{|X-\mu|\ge\varepsilon\}\le\frac{\sigma^2}{\varepsilon^2},
 \qquad\varepsilon>0.
\]
这里期望是按概率加权的平均值，方差衡量观测偏离平均值的程度，$\Pr$ 表示事件发生的概率。不等式说明：偏离平均值至少 $\varepsilon$ 的概率，有一个由方差决定的上界。
若 $X_1,\ldots,X_N$ 独立同分布，具有上述期望和方差，则样本均值 $\overline X_N$ 满足
\[
 \Pr\{|\overline X_N-\mu|\ge\varepsilon\}
 \le\frac{\sigma^2}{N\varepsilon^2}\longrightarrow0.
\]
这里趋于稳定的是\textbf{样本均值}，不是每条观测的噪声都消失了。切比雪夫不等式可以用于证明这一形式的大数定律，但不等式与大数定律本身不是同一个命题。

若样本只来自某类平台用户，就可能不能代表全部人群；若测量设备始终偏高，增加记录也不能消除这种系统偏差。模型还可能因时间、地区或行为习惯变化而失效，这称为分布变化。使用数据时，除了预测误差，还应考虑采集授权、隐私及对不同群体的影响。

\paragraph{要点归纳}
大数据思维用多维观测和统计关系改善不确定条件下的判断。数据量只有在带来有效信息时才有帮助；增加样本能减少某些随机误差，却不能自动纠正不具代表性的采样与系统偏差。相关关系可以成为预测线索，若要解释原因或采取干预，还需要进一步证据。

\section{大数据应用与定量决策}\label{sec:cases}
\subsection{推荐、搜索与预测}
推荐系统需要从已知偏好推测未知偏好。评分矩阵的行对应用户，列对应物品，元素记录评分；用户信息矩阵记录年龄、职业等属性，物品信息矩阵记录类型、内容等属性。后两者能够提供辅助信息，不同推荐方法对这些信息的依赖不同。

\textbf{评分缺失不等于评分为零。}缺失表示没有观测，零分则是一种已知评价。协同过滤利用不同用户或物品之间的共同结构进行预测，低秩矩阵补全是其中一种建模途径，见第\ref{sec:lowrank}节。预测后再按照偏好、候选范围和应用要求形成推荐。

搜索记录能够提供及时的行为信号，但搜索行为与目标现象之间的关系仍须检验。流感预测案例的基本过程是：选择可能相关的搜索特征，利用历史搜索量与实际疾病观测训练模型，再与新的实际观测比较。其作用是补充较及时的趋势信号；新闻关注和搜索习惯也会改变搜索量，因此不能把搜索次数直接当作患病人数。

广告匹配同样利用用户、内容和点击记录估计相关程度，购物关联分析则寻找共同出现的购买模式。药物用途探索从已有记录中寻找新的候选联系，再交由后续实验验证。这些应用共同体现了\textbf{从大量记录发现线索，再通过验证判断线索能否使用}的过程。

\subsection{多源数据与实时判断}
交通管理需要同时了解不同路段的状态。传统方法依赖固定观察、上报和汇总，覆盖与更新速度受到限制；移动终端产生的位置、速度和时间记录，使动态路况估计成为可能。数据的价值来自覆盖范围、实时性和相互关联，也取决于能否区分真实拥堵与停车、定位误差等其他情况。

饮食问答除了问题文字，还带有地域、时间和用户背景等维度。按这些维度组织记录，可以比较不同地区和群体的关注点；但提问者的关注不能直接等同于总体人群的实际饮食比例。选举等群体预测也要区分观测样本与目标人群，不能把网络发言简单当作全体人的选择。

异常检测则把个案与同类对象的统计规律比较。例如，用电模式异常可以作为进一步调查的线索，企业收入与同规模、同地区企业差异较大也可能需要核查。\textbf{异常不等于违规，统计筛查也不等于事实认定}；数据分析首先帮助缩小核查范围。

多源信息还可以相互补充。大庆油田的案例中，公开照片、新闻报道等分别提供间接线索，把这些线索放在一起，便能推断单条材料未直接说明的情况。这里的关键是比较不同来源能否相互印证，并区分已知事实与据此作出的推断。

\subsection{行为记录、社会关系与隐私}
用户行为不仅反映个人偏好，也可能受所处环境影响。锚定效应指先接触的数值或评价会影响后续判断；推荐中的评分因此未必是彼此独立、完全稳定的偏好测量。把这种行为机制纳入分析，有助于解释评分结构，并检验是否能改善推荐误差。

通话记录还能构成社会关系网络：用户是节点，联系是边，联系频率等可作为边的属性。网络联系的多样性可以与地区经济指标比较，通话模式也可能与年龄、性别或关系类型有关。这些案例把个体行为转成群体结构特征，展示了数据的多维价值；相关图形本身仍不能证明社会关系多样性直接导致经济增长。

跨域预测进一步说明，数据的潜在用途可能超出采集时的表面目的。把许多用户的点赞记录排成矩阵，再从中提取少量共同特征，就能用这些特征预测其他属性。课件采用奇异值分解提取特征，再用逻辑回归或线性回归作预测，后文将分别解释低秩表示与分类、回归的作用。这一过程也说明，看似普通的公开行为可能泄露未直接填写的私人信息。

评价时还应辨认指标。均方根误差记为
\[
 \operatorname{RMSE}=\sqrt{\frac1N\sum_{i=1}^N(\widehat b_i-b_i)^2}
\]
其中 $\widehat b_i$ 是预测值，$b_i$ 是实际值。RMSE 衡量数值预测的误差，通常越小越好。分类时，改变判定阈值，会同时改变\textbf{真实正类中被判为正类的比例}与\textbf{真实负类中被误判为正类的比例}。ROC 曲线描述这两者的权衡，AUC 是这条曲线下面积，反映模型区分正负样本的排序能力，并不等于固定阈值下的准确率。数据应用应同时考虑预测用途、验证结果与隐私影响。

\subsection{数据、模型与算法的配合}
机器翻译可以依靠人工编写词法、语法规则，也可以从大量对齐语料中学习语言之间的对应。后一种途径用数据和计算补充难以逐条编写的规则，说明复杂任务可以通过简单模型的组合或学习得到的表示来处理。

围棋程序则同时需要局面判断与行动搜索。训练数据帮助建立局面评价和落子选择的模型；蒙特卡罗树搜索通过对后续行动进行选择、扩展、模拟或评估，并回传结果，逐步把计算集中到更有希望的分支。模型提供判断，搜索组织决策过程，二者配合才形成完整系统。

这些案例表明，数据本身并不直接给出答案。还需要确定表示、模型、学习方法和计算过程，才能把记录转化为可使用的预测。
解释一个应用案例时，应把具体链条接完整：流感案例用搜索记录估计疾病趋势，推荐案例用已知偏好推测未知偏好，交通案例用动态位置和速度估计路况。每个例子都要继续检验观测是否可靠、关系能否延续，以及预测将怎样影响行动，不能只说“用了大数据，所以效果更好”。

\section{机器学习与深度学习}\label{sec:ml}
\subsection{人工智能与学习方法的层次}
人工智能关注智能任务的实现，机器学习则通过数据学习其中的规律。语音识别、图像识别、翻译和问答都属于人工智能的任务；完成这些任务可以采用知识规则，也可以采用学习方法。

知识库方法把事实和规则明确编码，再用逻辑规则推理。它便于表达清楚的知识，但对难以形式化描述的感知任务，逐条编写规则十分困难。机器学习从样本中学习模式，可以用于分类、预测和发现结构，但其效果仍受数据和模型限制。

图灵测试通过对话中能否区分机器与人来讨论机器智能。这是一种特定的行为评价方式；某项识别任务做得好，并不能据此证明机器具备人的全部智能。

几类方法之间具有如下关系：
\[
 \text{深度学习}\subset\text{表示学习}\subset\text{机器学习}\subset\text{人工智能}.
\]
表示学习让模型从数据中学习有效表示；深度学习进一步把简单表示逐层组合为复杂表示。并非所有机器学习都需要深层网络，也并非所有人工智能方法都通过数据训练。

\subsection{样本、特征与学习任务}
要让模型利用数据，需要先确定每个样本怎样表示，以及要预测什么。在监督学习中，记带标签的训练数据为
\[
 D=\{(\va_i,b_i):i=1,\ldots,N\},\qquad \va_i\in\R^d.
\]
其中 $\va_i$ 是第 $i$ 个样本的特征向量，$b_i$ 是相应标签或响应。模型参数与输入特征是不同对象：特征描述样本，参数则在训练中被调整，用来建立输入与输出之间的联系。
例如猫狗分类中，一张图片是一条样本，各像素的数值可以作为输入表示，猫或狗是标签；网络的连接权重是参数。像素来自图片，权重需要由训练估计，不能把二者当成同一组数字。
\begin{itemize}
\item \textbf{监督学习：}利用带标签样本学习输入到输出的关系。输出为连续数值时称为回归，输出为类别时称为分类。
\item \textbf{无监督学习：}样本没有给定目标标签，任务是发现聚类、低维表示等结构。聚类得到的分组不同于已知类别下的分类。
\item \textbf{强化学习：}通过与环境互动，利用动作产生的奖励学习策略，关注长期行动效果。奖励不同于每个样本的正确类别标签。
\end{itemize}

特征选择从已有特征中选取子集；特征提取通过变换构造新的特征；表示学习让这种变换也由数据学习。例如，手工设计图像边缘特征属于特征提取，网络逐层学到边缘、局部形状和更抽象结构则体现表示学习。特征应服务于目标，维度多并不必然优于维度少。

表示方式还会改变分类的难易。考虑同心分布的两类平面点：内层点云被外层环状点云围住时，在原来的笛卡尔坐标中不能用一条直线将它们分开。若改用极坐标，取点到共同中心的距离 $r$ 作为特征，就可用一个半径阈值区分内外两类。改变的是表示及其下的分界形式，原来的样本和类别并没有改变。

\subsection{训练、损失与泛化}
训练通过调整参数减小损失，泛化则要求模型在未用于训练的数据上仍有较好表现。设 $f(\va;\vx)$ 为参数 $\vx$ 下的预测，$\ell$ 为单个样本的损失，常见训练模型为
\begin{equation}\label{eq:learning}
 \boxed{\min_{\vx\in\mathcal W}\frac1N\sum_{i=1}^N
 \ell\bigl(f(\va_i;\vx),b_i\bigr)+\lambda R(\vx).}
\end{equation}
第一项是经验损失，衡量训练数据上的拟合程度；$R(\vx)$ 是正则项，表达对参数大小、稀疏性等结构的要求；$\lambda\ge0$ 控制两者的相对权重，$\mathcal W$ 是允许的参数集合。将损失之和改为平均值时，正则参数也应相应调整，才能保持同一个优化问题。
从计算过程读这个式子，就是先用当前参数给每条样本作预测，再将预测与已知标签比较，计算损失；把所有样本损失平均以后，再加上对参数的额外要求。训练寻找使这个总目标较小的参数。式中的“经验”指来自现有样本，“正则”则是除拟合数据以外的选解要求，第\ref{sec:regression}节将用具体例子解释。

在总体分布上希望减小的是期望损失
\[
 \E\bigl[\ell(f(\va;\vx),b)\bigr].
\]
经验损失由有限样本近似这一期望，因而训练误差小不自动保证总体误差小。

\textbf{欠拟合}指模型未能充分描述数据中的主要结构；\textbf{过拟合}指模型过多适应训练样本中的偶然细节，导致新数据上的表现较差。增加模型复杂性可能缓解欠拟合，却也可能增加过拟合风险。正则化改变选解标准，优化算法则负责求解给定的训练问题；把训练目标降得更低，并不总能改善泛化。

训练集用于估计参数，验证集用于选择模型和正则参数等设置，测试集用于评价最终方案。交叉验证把数据轮流划分为训练部分和验证部分，以比较候选设置；反复根据测试结果调参，会使测试集失去独立评价作用。对时间序列，还应保持时间先后，避免把未来信息用于过去的预测。

\subsection{神经网络的组成}
一个神经元先把输入乘以相应权重并相加，再加偏置，最后经过激活函数。例如两个输入时，先算 $z=w_1a_1+w_2a_2+c$，再算 $t(z)$。权重决定各输入的作用，偏置调整整体位置，激活函数决定怎样把加权结果变为输出。许多这样的单元连接起来，便构成神经网络。

多层感知机由输入层、若干隐藏层和输出层组成，隐藏层就是输入与输出之间的中间层，信息沿前向连接传播。记第 $l$ 层的输出为 $\boldsymbol h^{(l)}$，则一层运算可写为
\[
 \boldsymbol h^{(0)}=\va,\qquad
 \boldsymbol z^{(l)}=W^{(l)}\boldsymbol h^{(l-1)}+\boldsymbol c^{(l)},\qquad
 \boldsymbol h^{(l)}=t_l(\boldsymbol z^{(l)}).
\]
$W^{(l)}$ 是连接权重，$\boldsymbol c^{(l)}$ 是偏置，$t_l$ 是激活函数，通常逐分量作用。所有层的权重和偏置共同构成训练参数。

常见激活函数为
\[
 \operatorname{sigmoid}(z)=\frac1{1+e^{-z}},\qquad
 \operatorname{ReLU}(z)=\max\{0,z\},\qquad
 \operatorname{Heaviside}(z)=
 \begin{cases}1,&z\ge0,\\0,&z<0.\end{cases}
\]
sigmoid 把实数压入 $0$ 与 $1$ 之间；ReLU 保留正数，把负数变为零；Heaviside 按零这个阈值输出 $0$ 或 $1$。非线性激活使网络能够表达非线性关系。若每层都只做加权相加再加偏置，即仿射变换，那么多层复合仍是一次仿射变换，增加层数并未获得非线性表达能力。输出层的形式还要与回归或分类任务相适应。

\subsection{不同网络处理不同结构}
图像具有局部空间结构，序列则包含前后依赖，因此网络结构应与数据结构相适应。

\textbf{卷积神经网络（CNN）}用局部连接和权重共享提取空间特征。一个卷积核在图像上滑动，对局部区域与核的对应元素相乘求和，得到特征图。若图像为 $I$、核为 $K$，一种常用计算形式为
\[
 S_{ij}=\sum_{u=1}^k\sum_{v=1}^k K_{uv}I_{i+u-1,j+v-1}.
\]
输出尺寸由核大小、步长和边界填充共同决定。与全连接相比，局部连接减少连接数，权重共享使同一种局部模式能在不同位置被识别；池化可进一步汇总局部特征。严格按信号处理术语，上式为互相关，神经网络中通常也称为卷积。

\textbf{转置卷积}使用与某个卷积线性映射对应的转置映射，常用于增加空间分辨率和生成图像。它并不是一般意义上恢复原图的逆卷积，也不能保证把丢失信息自动还原。

\textbf{循环神经网络（RNN）}通过隐藏状态保留序列上下文，例如
\[
 \boldsymbol h_t=t\bigl(W_h\boldsymbol h_{t-1}+W_a\va_t+\boldsymbol c\bigr).
\]
其中 $W_h,W_a$ 分别是作用于前一状态和当前输入的权重，$\boldsymbol c$ 是偏置，$t$ 是激活函数。当前状态因而同时依赖当前输入与前一状态。各时刻使用同一组权重，适合处理语音、文字等序列。状态把有关历史压入一组数值，并不是显式存储全部历史。

语音识别可以把观测声音看作待解码的信号，再结合语言规律选择较可能的文字。基于资料的问答则需要先限定可回答范围，检索候选片段，再判断相关性并组织回答。二者都说明：先确定任务和表示，才能选择合适的模型与学习过程。

\paragraph{要点归纳}
机器学习利用样本估计规律，表示学习还让输入的有效表示由数据学得，深度学习通过多层复合建立这种表示。模型结构要与空间、时间等数据关系相配合。训练损失反映现有样本上的拟合，泛化要求新样本上也有效，因此训练、验证与测试承担不同职责，不能互相代替。

\section{最优化问题与建模方法}\label{sec:model}
\subsection{推断与决策}
学习模型给出预测，行动选择还需要考虑资源与目标。例如，预测某条道路会拥堵以后，仍要比较可行路线的时间和成本；预测用户偏好以后，还要决定展示哪些物品。
优化既可以用来选择行动，也可以用来估计模型参数。投资中要选择的是资金比例，训练中要选择的是使损失较小的参数；两类问题的未知量不同，却都需要明确变量、目标与约束。

\textbf{数据科学}主要从已知记录推断未知信息，涉及收集、清洗、表示、建模和评价。\textbf{运筹学}主要研究在给定目标和约束下怎样决策，通过数学模型比较行动方案。二者相互支持：预测为决策提供参数，优化用于估计学习模型，学习也可用于近似复杂决策中的价值或策略。

俾斯麦海战例把侦察方向与航线选择看成双方策略。一般地，设一方希望增大收益，对方希望压低这一收益。若 $a_{ij}$ 表示前者选择策略 $i$、后者选择策略 $j$ 时前者的收益，则保守决策先比较各策略的最坏收益，再取其中最大的一个：
\[
 \max_i\min_j a_{ij}.
\]
希望压低收益的一方则相应比较 $\min_j\max_i a_{ij}$。当某个位置同时是所在行的最小值、所在列的最大值时，形成纯策略鞍点；一般情况下二者未必相等。这个模型强调\textbf{在对方也会行动的条件下比较策略}，不能只挑某一种有利情形下收益最大的方案。

反潜与护航案例则从损失、搜索效果等观测出发，改进编队、搜索和资源使用。它与博弈案例共同说明：决策需要明确目标、可选行动、限制条件和比较标准。

\subsection{决策变量、目标函数与约束}
建立优化模型，就是确定可以改变的量、比较方案的标准和必须满足的条件。一般最小化问题写为
\begin{equation}\label{eq:optimization}
 \boxed{\min_{\vx\in\mathcal X} f(\vx).}
\end{equation}
\begin{definition}
$\vx$ 称为\textbf{决策变量}，$f$ 称为\textbf{目标函数}，$\mathcal X$ 称为\textbf{可行域}。属于 $\mathcal X$ 的点称为可行点；若 $\vx^*\in\mathcal X$ 且对所有 $\vx\in\mathcal X$ 有
\[
 f(\vx^*)\le f(\vx),
\]
则 $\vx^*$ 为全局最优解，$f(\vx^*)$ 为最优值。
\end{definition}
可行域常由等式与不等式组成，例如
\[
 \mathcal X=\{\vx:g_i(\vx)\le0,\ i=1,\ldots,p;
 \ h_j(\vx)=0,\ j=1,\ldots,q\}.
\]
无约束问题取 $\mathcal X=\R^n$。最大化 $f$ 可以改写为最小化 $-f$，两者的最优解相同，最优值互为相反数。
后文公式中的 $\text{s.t.}$ 是“满足以下约束”的缩写，列在它后面的条件必须同时满足。

最优解不一定存在。例如 $\min_{x\in\R}e^x$ 的下确界是 $0$，却没有任何有限 $x$ 取得这个值。可行、达到最优与只有下确界，是不同的情况；目标函数也可能向下无界。

\begin{example}[投资组合]
设 $x_i$ 为投在第 $i$ 项资产的资金比例，$\mu_i$ 为预期回报率，$\Sigma$ 为回报率的协方差矩阵。以方差度量风险，在不允许卖空且要求预期回报至少为 $r_0$ 时，可写为
\[
 \begin{aligned}
 \min_{\vx}\quad&\frac12\vx^{\mathsf T}\Sigma\vx,\\
 \text{s.t.}\quad&\boldsymbol\mu^{\mathsf T}\vx\ge r_0,\qquad
 \sum_{i=1}^n x_i=1,\qquad x_i\ge0.
 \end{aligned}
\]
\end{example}
资金比例是变量，风险是目标。协方差矩阵 $\Sigma$ 记录各资产回报自身及彼此的共同波动，$\vx^{\mathsf T}\Sigma\vx$ 就是按比例组合后的总回报方差。$\boldsymbol\mu^{\mathsf T}\vx$ 是各资产预期回报按投资比例加权后的总预期回报；$\sum_i x_i=1$ 表示全部资金都被分配；$x_i\ge0$ 表示不允许出现负的持有比例，即不允许卖空。这三项要求各有不同作用。目标前的 $1/2$ 不改变最优投资方案，只是方便计算。

风险还可用风险价值（VaR）、条件风险价值（CVaR）等指标衡量。给定置信水平 $\alpha$，VaR 是损失的相应分位数，可理解为有 $\alpha$ 比例的损失不超过的门槛；在连续分布情形下，CVaR 是超过这一门槛后的尾部平均损失。这是在描述损失分布，不是在保证投资回报。不同指标可能使模型具有不同的光滑性与计算性质，因此实际目标相似，不意味着数学模型和求解算法相同。

\subsection{模型的类别与性质}
同一个实际问题可以有不同模型，而模型的性质会影响求解难度。
\begin{itemize}
\item \textbf{线性规划：}目标为线性函数，等式、不等式约束也是线性的。常数项不影响这一分类。
\item \textbf{二次规划：}目标为二次函数，约束为线性等式或不等式。投资组合的方差模型属于此类。
\item \textbf{非线性规划：}目标或约束包含非线性关系。二次规划也是非线性规划的一类。
\item \textbf{整数与组合优化：}部分或全部变量只能取整数，或表示离散选择。
\item \textbf{非光滑优化：}目标或约束含不可微结构，如绝对值在零点不可微。无导数方法则指不直接利用导数信息的方法，所处理的函数不一定不可微。
\item \textbf{半定规划：}在线性等式、不等式及矩阵半正定约束下优化线性目标。对称矩阵 $X\succeq0$ 表示对任意向量 $\boldsymbol v$，都有 $\boldsymbol v^{\mathsf T}X\boldsymbol v\ge0$。
\end{itemize}
稀疏优化与低秩优化按希望得到的解结构命名；随机优化考虑随机变量或期望目标，鲁棒优化考虑不确定条件下的性能。这些分类可以交叉，不能看成互不相容的几类。

\begin{definition}[凸性]
若对 $\vx,\boldsymbol y\in\mathcal X$ 和 $t\in[0,1]$，都有 $t\vx+(1-t)\boldsymbol y\in\mathcal X$，则 $\mathcal X$ 为凸集。若在凸集上
\[
 f(t\vx+(1-t)\boldsymbol y)\le t f(\vx)+(1-t)f(\boldsymbol y),
\]
则 $f$ 为凸函数。在凸集上最小化凸函数，称为凸优化问题。
\end{definition}
凸集要求任取两个允许方案，它们之间的线段仍然允许。例如区间 $[0,1]$ 是凸集，只有 $0,1$ 两点的集合不是，因为中点 $1/2$ 不在其中。凸函数的不等式则表示：两点间的函数图像不高于连接两端的弦线；它为排除“局部较好却远处还有更好”的情形提供依据。
凸优化的重要性质是局部最优解也是全局最优解，但凸性本身不保证解一定存在或唯一。非凸问题可能有不同的局部最优解，算法对初始点的依赖通常更值得关注。

\subsection{目标与约束的设计}
观测通常含有误差，目标函数需要表达怎样衡量误差，以及希望解具有什么结构。设预测为 $\phi_i(\vx)$、观测为 $b_i$，残差为 $r_i(\vx)=\phi_i(\vx)-b_i$。三种常见选择为
\[
 \min_{\vx}\sum_i r_i(\vx)^2,\qquad
 \min_{\vx}\sum_i|r_i(\vx)|,\qquad
 \min_{\vx}\max_i|r_i(\vx)|.
\]
平方损失对大残差给予更大的惩罚；绝对值损失相对减弱极端残差的影响；最大残差目标直接控制最坏误差。选择哪一种，应由误差特点和应用要求决定。

成本、损失和收益也可直接作为目标。若变量是函数，目标可能是定义在函数空间上的泛函，例如能量；通过离散化，可把某些问题转成有限维优化问题。约束则表达资源、非负性、守恒关系、正交关系或其他物理要求。模型合理性不仅表现为拟合已知数据，还要看能否反映实际限制，并与已有方法作公平比较。

\subsection{等价变换}
有时模型难以直接求解，可以先改变表达形式，同时保留原问题的含义。例如
\[
 \min_{\vx} f(A\vx+\vb)
 \quad\Longleftrightarrow\quad
 \min_{\vx,\boldsymbol y} f(\boldsymbol y),\quad
 \text{s.t. }\boldsymbol y=A\vx+\vb.
\]
右边引入了新变量，但没有放宽 $\boldsymbol y$ 与 $\vx$ 的关系。同样，$\min_{\vx}\{h(\vx)+r(\vx)\}$ 可拆为
\[
 \min_{\vx,\boldsymbol y}h(\vx)+r(\boldsymbol y),
 \qquad\text{s.t. }\vx=\boldsymbol y.
\]
对不等式 $g(\vx)\le0$，引入非负变量 $s$，可等价写为 $g(\vx)+s=0,\ s\ge0$。这个 $s$ 虽称为松弛变量，但该改写\textbf{没有扩大原变量的可行域}。

引入目标值的上界变量 $t$，$\min_{\vx} f(\vx)$ 还可等价写为
\[
 \min_{\vx,t}t,\qquad\text{s.t. }f(\vx)\le t.
\]
等价变换通过消去辅助变量可回到原问题；它改变计算组织方式，而不改变原变量下的可行方案和最优目标。

\subsection{替代目标与约束松弛}
另一类处理会改变原问题。例如，用 $\ell_1$ 范数替代非零分量计数，或用核范数替代矩阵秩，保留了促进稀疏或低秩的意图，却一般不能保证解完全相同。

约束松弛则把可行域 $\mathcal X$ 扩大为 $\mathcal X_R\supseteq\mathcal X$。例如用 $0\le x\le1$ 替代 $x\in\{0,1\}$，或把等式约束放宽为单侧不等式。对同一最小化目标有
\[
 \inf_{\vx\in\mathcal X_R}f(\vx)\le
 \inf_{\vx\in\mathcal X}f(\vx).
\]
其中 $\inf$ 表示下确界，即不超过任何目标值的那些数中最大的一个；即使没有某个可行点取得它，也可以讨论这个下界。松弛后的下确界不高于原问题，因而提供下界，但解可能不满足原约束。若采用下界函数 $f_R\le f$ 替代目标，也能形成相应下界；并非所有替代目标都自动具有这一性质。

单位球面约束要求向量长度为 $1$。把向量分量的两两乘积排成矩阵 $X=\vx\vx^{\mathsf T}$，就将向量问题提升到矩阵变量。这样的矩阵恰好满足
\[
 X\succeq0,\qquad\operatorname{tr}(X)=1,\qquad\rank X=1.
\]
这里 $\operatorname{tr}(X)$ 是对角元素之和，等于 $\norm{\vx}_2^2$。去掉秩为 $1$ 的约束，就得到更大的集合 $X\succeq0,\ \operatorname{tr}(X)=1$。若结果仍为秩 $1$，可以分解回单位向量；若秩大于 $1$，则不能直接写成单个 $\vx\vx^{\mathsf T}$。在具体模型中，还须把目标与其他约束相应表达在提升变量下，才构成完整松弛。

\paragraph{要点归纳}
建模先明确哪些量可调整，再把好坏标准写成目标，把预算、物理条件等写成约束。等价变换增加辅助变量但保留原变量的方案；约束松弛放宽允许范围；用易处理目标替代困难目标则改变了比较标准。这三件事作用不同，不能只因形式更简单就称为完全等价。

\section{回归、似然与正则化}\label{sec:regression}
\subsection{从预测式读懂回归}
回归要根据已知特征预测一个连续数值。例如，用房屋面积和房龄预测房价，就是用两个输入预测一个数值输出。先不考虑具体单位，把这样的线性关系写为
\[
 \widehat b_i=x_1a_{i1}+x_2a_{i2}+c.
\]
$i$ 表示第几条样本，$a_{i1},a_{i2}$ 是已知的两个特征，$b_i$ 是已知观测值，$\widehat b_i$ 是预测值。系数 $x_1,x_2$ 和截距 $c$ 是待估参数；$c$ 给出两个输入都为零时的基准预测。训练要调整的是这些参数，而不是把已观测的数据任意改成想要的数值。

\begin{example}[预测与残差]
在一个只说明计算关系的例子中，取输入 $(a_{i1},a_{i2})=(2,1)$，参数 $(x_1,x_2,c)=(3,-1,1)$，则预测值为 $3\times2-1\times1+1=6$。若实际观测为 $b_i=7$，则残差“预测减观测”为 $6-7=-1$，残差平方为 $1$。
\end{example}
一条样本只有一个残差，全部样本各有各的残差。\textbf{最小二乘}就是选择一组参数，让这些残差的平方加起来尽量小：
\[
 \min_{x_1,x_2,c}\frac12\sum_{i=1}^N
 (x_1a_{i1}+x_2a_{i2}+c-b_i)^2.
\]
平方使正负误差不会抵消，也给较大的误差更重的惩罚。前面的 $1/2$ 只是方便求导，不改变最优参数。

将每条样本的特征放成矩阵 $A$ 的一行，将输出排成向量 $\vb$，将待估参数排成列向量 $\vx$，同一件事可以缩写为
\[
 \boxed{\min_{\vx}\frac12\norm{A\vx-\vb}_2^2.}
\]
这里 $A\vx$ 是全部样本的预测，$A\vx-\vb$ 是全部残差，$\norm{A\vx-\vb}_2^2$ 就是残差平方和。若保留截距，在 $A$ 中增加一列 $1$，并把 $c$ 放入参数向量即可。\textbf{矩阵写法没有增加新的任务，只是一次写出全部样本。}

“线性回归”的线性指对参数线性；“最小二乘”指采用平方和损失。非线性预测函数同样可以使用平方损失，得到非线性最小二乘。观测有误差时，通常不要求每条预测都恰好等于观测，而是比较哪组参数的总误差较小；解也不一定唯一。

\subsection{最大似然为什么得到最小二乘}
误差度量还可以从概率假设得到。设已经观察到一组数据，尝试不同参数时，模型对这组数据给出的概率或密度会不同。\textbf{最大似然估计}选择让已经观测到的数据在模型下更可能出现的参数。

若样本独立，同一参数下整组样本的联合概率或密度是各项的乘积。记第 $i$ 条观测为 $z_i$、模型为 $p(z_i;\vx)$，则
\[
 L(\vx)=\prod_{i=1}^N p(z_i;\vx).
\]
数据在这里已固定，比较的是不同参数 $\vx$。这不是计算“参数本身有多大概率”。对于连续观测，使用的是概率密度，它描述概率在数值附近的集中程度，单点密度值并不是该精确数值出现的概率。

乘积不便处理，可以取对数把它变成求和。对数严格递增，不改变最大值的位置；再取负号，就把最大化改为最小化：
\[
 \max L(\vx)\quad\Longleftrightarrow\quad
 \min -\sum_i\log p(z_i;\vx).
\]
在线性回归中，若误差相互独立、均值为零，并服从共同方差 $\sigma^2>0$ 的高斯分布，则负对数似然为
\[
 \text{与参数无关的常数}
 +\frac1{2\sigma^2}\sum_i(\widehat b_i-b_i)^2.
\]
常数不影响选哪组参数，正的公共倍数也不影响最小值的位置，所以此时\textbf{最大似然选出的回归参数，恰好就是最小二乘参数}。求这组参数不必先知道共同方差的具体值。若误差独立服从共同尺度的零中心 Laplace 分布，则对应的是绝对残差和，而不是残差平方和。因此损失的选择应与误差特点相配合。

\subsection{为什么还要控制系数}\label{sub:regularization}
训练误差小，仍可能没有把模型确定好。考虑两列特征在训练数据中总相等的情形：预测式为 $x_1a_{i1}+x_2a_{i2}$，只要 $x_1+x_2=1$，预测都等于这一条样本的共同特征值。

于是系数 $(100,-99)$ 与 $(0.5,0.5)$ 在这些训练样本上给出完全相同的预测。但是，对输入 $(1.01,1)$，前者预测为 $2$，后者预测为 $1.005$。第一组系数依靠两个大数相互抵消，输入发生很小的差异，抵消就被破坏。这个例子说明，\textbf{仅比较训练误差，可能无法排除对小变化敏感的参数}；并不是说任何较大的系数都错误。

因此，除了拟合观测，还可以对参数提出额外的选择标准，例如抑制过大的系数，或希望只使用少数特征。这种做法称为\textbf{正则化}。它在目标中体现为
\[
 \boxed{\text{总目标}=\text{拟合误差}+\lambda\times\text{参数惩罚}.}
\]
误差项要求预测接近观测，惩罚项表达希望参数具有什么性质，$\lambda\ge0$ 控制两者的相对权重。正则化可能牺牲一点训练拟合，换取较稳定或较简单的模型；能否改善新数据上的表现，仍须用验证数据检验。

\subsection{岭回归与LASSO：缩小和筛选}
下面省去单独的截距，只对特征系数 $\vx=(x_1,\ldots,x_n)^{\mathsf T}$ 加惩罚。\textbf{岭回归}采用系数平方和，是 Tikhonov 正则化的一种形式：
\begin{equation}\label{eq:ridge}
 \boxed{\min_{\vx}\underbrace{\frac12\norm{A\vx-\vb}_2^2}_{\text{拟合误差}}
 +\lambda\underbrace{\sum_jx_j^2}_{\text{系数平方和}},\qquad\lambda>0.}
\end{equation}
较大的系数会付出较重的代价，因此岭回归倾向于把系数缩小，改善不稳定估计。在全部特征系数都受到这一惩罚时，线性平方损失模型有唯一解。但系数缩小，不等于系数恰好为零。

\textbf{LASSO}把惩罚改为系数绝对值之和：
\begin{equation}\label{eq:lasso}
 \boxed{\min_{\vx}\underbrace{\frac12\norm{A\vx-\vb}_2^2}_{\text{拟合误差}}
 +\lambda\underbrace{\sum_j|x_j|}_{\text{系数绝对值之和}}.}
\end{equation}
LASSO 除了缩小系数，还会促进部分系数恰好为零。\textbf{稀疏}就是指多数分量为零、只有少数非零。在预测式中，一个特征的系数为零，表示模型不再使用这个特征，因此 LASSO 可以用于特征选择。

\begin{example}[缩小与变为零的区别]
用一个参数 $t$ 看清两种惩罚。设不加惩罚时，拟合损失为 $\frac12(t-0.2)^2$，它在 $t=0.2$ 时最小。
\begin{enumerate}
\item 加平方惩罚后，求解 $\min_t\{\frac12(t-0.2)^2+\lambda t^2\}$。令导数为零，得 $(t-0.2)+2\lambda t=0$，故最优参数为 $t=0.2/(1+2\lambda)$。
\item 加绝对值惩罚后，求解 $\min_t\{\frac12(t-0.2)^2+\lambda|t|\}$。在 $t>0$ 时，导数为 $t-0.2+\lambda$；负数的拟合误差大于其对应的正数，而惩罚相同，因而不必选择负数。还需比较边界 $t=0$，最后得到最优参数 $t=\max\{0.2-\lambda,0\}$。
\end{enumerate}
\begin{center}
\begin{tabular}{ccc}
\toprule
$\lambda$ & 平方惩罚后的参数 & 绝对值惩罚后的参数\\
\midrule
$0$ & $0.2$ & $0.2$\\
$0.1$ & 约 $0.167$ & $0.1$\\
$0.3$ & $0.125$ & $0$\\
\bottomrule
\end{tabular}
\end{center}
绝对值惩罚使参数一旦不够“值得保留”，就可以直接停在零；平方惩罚在这个例子中使非零参数逐渐缩小。这解释了两者作用的差别，但并不意味着一般 LASSO 都能逐个系数按此式求解。
\end{example}

$\lambda$ 太小，额外惩罚的作用弱，可能仍过拟合或不稳定；太大则过分限制参数，可能连有用的关系也表达不好，造成欠拟合。通常用验证集选择 $\lambda$，而不只比较训练误差。

系数大小还受单位影响：同一长度改用米或厘米表示，对应系数就会改变，而原来的预测可以不变。因此比较惩罚前通常要统一尺度，例如将每个特征减去其均值，再除以相应标准差，这称为标准化；常量特征不能直接除以零标准差，应另行处理。截距常单独处理而不加同样惩罚。$\ell_1$ 促进稀疏，也不保证任意数据、任意参数下都能正确筛出真实相关特征。

\paragraph{要点归纳}
岭回归与 LASSO 都在拟合之外增加参数惩罚。岭回归使用系数平方和，主要缩小系数、改善稳定性；LASSO 使用绝对值之和，促进部分系数为零，从而选择特征。正则化改变的是选解标准，优化算法负责求解这个标准下的问题，验证则判断它对新数据是否有效。

\subsection{约束形式与结构化正则}
正则化也可以直接规定预算。例如 $\norm{\vx}_1=\sum_j|x_j|$，模型
\[
 \min_{\vx}\frac12\norm{A\vx-\vb}_2^2,
 \qquad\text{s.t. }\norm{\vx}_1\le\tau
\]
表示“在系数绝对值总预算 $\tau$ 内，把拟合误差降到最低”。另一种写法是
\[
 \min_{\vx}\norm{\vx}_1,
 \qquad\text{s.t. }\norm{A\vx-\vb}_2\le\delta,
\]
即“允许残差不超过 $\delta$，在此条件下使系数绝对值之和最小”。岭回归也可规定 $\norm{\vx}_2\le\tau$ 的预算，其中 $\norm{\vx}_2=\sqrt{\sum_jx_j^2}$。这些形式表达拟合与结构之间的权衡，在适当条件和参数对应下可以有相同最优解，但不能随意把 $\lambda,\tau,\delta$ 当作同一个数。

损失项和正则项的作用也应区分。将平方残差改为 $\norm{A\vx-\vb}_2$ 或 $\norm{A\vx-\vb}_1$，改变的是衡量拟合误差的方式；将正则项改为不同形式，改变的是希望参数具有什么结构。

若希望选择整组变量，或希望信号只在少数位置改变，则需要\textbf{结构化正则}。记 $I_g$ 为第 $g$ 组的变量指标，$n_g$ 为该组大小，$F$ 为一个给定的线性变换，各种要求可以概括为：
\begin{center}
\begin{tabular}{p{0.21\textwidth}p{0.35\textwidth}p{0.30\textwidth}}
\toprule
方法 & 加入的惩罚 & 希望得到的结构\\
\midrule
分组 LASSO & $\lambda\sum_g\sqrt{n_g}\norm{\vx_{I_g}}_2$ & 某些组整体为零，舍去整组特征。\\
稀疏分组 LASSO & 分组惩罚再加 $\lambda_2\norm{\vx}_1$ & 既选择组，也选择组内的特征。\\
变换域稀疏 & $\lambda\norm{F\vx}_1$ & 变换后的分量多数为零，原信号未必稀疏。\\
融合 LASSO & $\lambda_1\norm{\vx}_1+\lambda_2\sum_{j=2}^n|x_j-x_{j-1}|$ & 系数可为零，相邻系数也倾向于相同。\\
\bottomrule
\end{tabular}
\end{center}
分组惩罚先计算每组系数的长度，再把组长度加起来；组被选中不表示其内部每个系数都非零。若每组只有一个变量，分组 LASSO 退化为普通 LASSO。差分惩罚则把“相邻两项不同”计入代价，适合描述大段平稳、少量位置发生变化的信号。共同原则是：\textbf{希望促进哪些量为零，就围绕这些量选择合适的惩罚。}

\section{稀疏优化与低秩恢复}\label{sec:sparse}
\subsection{从一条方程的多个解理解稀疏}
方程可以限制未知数，却未必唯一决定未知数。考虑
\[
 x_1+2x_2=2.
\]
$(2,0)$、$(0,1)$、$(0.4,0.8)$ 都满足方程。若只检查是否满足观测，三者都可行。要选其中一个，就需要额外标准。

对一个向量 $\vx$，$\norm{\vx}_0$ 数非零分量的个数，$\norm{\vx}_1$ 加各分量的绝对值，$\norm{\vx}_2^2$ 加各分量的平方。前者常称 $\ell_0$“范数”，但严格说不是范数；$\ell_2$ 范数是平方和再开平方，最小化它与最小化其平方会选到相同的点。
\begin{center}
\begin{tabular}{cccc}
\toprule
可行解 & 非零分量数 $\norm{\vx}_0$ & 绝对值和 $\norm{\vx}_1$ & 平方和 $\norm{\vx}_2^2$\\
\midrule
$(2,0)$ & $1$ & $2$ & $4$\\
$(0,1)$ & $1$ & $1$ & $1$\\
$(0.4,0.8)$ & $2$ & $1.2$ & $0.8$\\
\bottomrule
\end{tabular}
\end{center}
若要求非零项最少，前两组都是最优解。若要求绝对值和最小，则选 $(0,1)$：由 $2=|x_1+2x_2|\le2(|x_1|+|x_2|)$，绝对值和至少为 $1$，达到等号还要求 $x_1=0$，故解唯一。若要求平方和最小，把 $x_1=2-2x_2$ 代入，有
\[
 x_1^2+x_2^2=5(x_2-0.8)^2+0.8,
\]
所以选 $(0.4,0.8)$。\textbf{三种标准选出的结果不同，不能把某一种标准的最优解当作其他标准的最优解。}

\subsection{稀疏恢复与压缩感知}
如果已知真实信号只有少数非零分量，就可以用这种稀疏性作为额外依据。在方程 $A\vx=\vb$ 相容且未知数多于独立方程时，通常有无穷多个解，直接寻找最稀疏者为
\[
 \min_{\vx}\norm{\vx}_0,\qquad\text{s.t. }A\vx=\vb.
\]
这里已知的是测量方式 $A$ 与测量结果 $\vb$，未知的是信号 $\vx$。约束要求信号符合测量，目标要求使用尽量少的非零分量。

直接优化非零个数一般很困难，常用绝对值和替代，得到\textbf{基追踪}：
\[
 \boxed{\min_{\vx}\norm{\vx}_1,\qquad\text{s.t. }A\vx=\vb.}
\]
在测量矩阵、稀疏程度等满足适当条件时，基追踪能够恢复真实稀疏信号。\textbf{压缩感知}利用这种结构，从较少测量恢复较高维信号；并不是任意少量观测都能还原任意信息。例如，若前一例的真实信号是 $(2,0)$，绝对值和最小化却会选择 $(0,1)$，并未恢复真实信号。两者都符合测量且都稀疏，仍不足以保证选对；$\ell_0$ 和 $\ell_1$ 也不能一般地宣布等价。

基追踪是凸问题，还可通过 $-t_j\le x_j\le t_j$ 把绝对值和写成线性规划。但当 $A$ 很大、稠密，或只能计算 $A\vx$ 与 $A^{\mathsf T}\boldsymbol y$ 而不能存储全部元素时，仍需选择能利用这些结构的算法。凸性有助于求解，并不消除全部计算成本。

最小 $\ell_2$ 范数选择欧氏长度最小的解，通常不会特意留下大量零分量。它本身也是合理标准，只是与稀疏恢复的任务不同。带噪观测则常改用残差容差或 LASSO 的惩罚形式，避免要求精确符合每一个含噪记录。

\subsection{推荐为什么能够利用低秩}\label{sec:lowrank}
一个用户没看过某部电影，评分表中的空白表示未知，而不是零分。考虑一个简化的两用户、两电影评分表：
\[
 \begin{array}{c|cc}
 &\text{电影甲}&\text{电影乙}\\\hline
 \text{用户甲}&5&?\\
 \text{用户乙}&5&3
 \end{array}
\]
仅凭这些数字，不能确定问号是多少。若额外假定两个用户的整行评分完全相同，则问号应为 $3$。用线代语言表达，这个 $2\times2$ 例子若要求矩阵秩为 $1$，两行须成比例；由于第一列同为非零数 $5$，比例只能为 $1$，所以第二列也必须相同。\textbf{补全来自观测与结构假设的结合，而不是空白本身给出了答案。}

矩阵的秩是线性无关行或列的最大个数。低秩意味着大量行列可以由少数共同模式表达；在推荐中，这相当于尝试用少数潜在偏好因素解释许多用户和电影。\textbf{观测稀疏}是已知位置少，\textbf{低秩}是行列之间的结构简单，两者不是同一性质。

一般地，已观测评分记作 $M_{ij}$，其位置集合为 $\Omega$，完整评分矩阵 $X$ 待求，则
\[
 \boxed{\min_X\rank X,\qquad
 \text{s.t. }X_{ij}=M_{ij},\quad(i,j)\in\Omega.}
\]
读成中文就是：\textbf{在保留已知评分的所有补全中，选择秩尽量低的矩阵。}低秩假设不一定符合实际，也不单独保证补全唯一。

\subsection{为什么把秩换成核范数}
秩数的是独立模式的个数，直接优化这种计数很困难。矩阵的奇异值是一组由矩阵分解得到的非负数，用来描述这些模式的大小；它们不是矩阵各位置上的评分。设这些数为 $\sigma_j$，则
\[
 \rank X=\#\{j:\sigma_j>0\},\qquad
 \norm{X}_*=\sum_j\sigma_j.
\]
例如奇异值为 $4,1,0$ 时，秩为 $2$，核范数为 $5$；把它们改为 $4,0.2,0$，秩仍是 $2$，但核范数降为 $4.2$。计数只看是不是零，求和还区分大小。一般矩阵的秩应由非零奇异值计数，不能直接改用非零特征值个数。

\textbf{核范数}就是奇异值之和。用它替代秩，得到
\[
 \boxed{\min_X\norm{X}_*,\qquad
 \text{s.t. }X_{ij}=M_{ij},\quad(i,j)\in\Omega.}
\]
这是凸问题，借助核范数促进部分奇异值为零，也就是促进低秩。它与向量中用 $\ell_1$ 替代 $\ell_0$ 的思想相似；替代是否能恢复真实矩阵，仍依赖观测数量、位置和矩阵结构，不能保证总与秩最小化同解。

\subsection{含噪评分与低秩因子}
若观测评分本身有误差，就不必强求全部已知位置精确相等，可以改为
\[
 \min_X\underbrace{\frac12\sum_{(i,j)\in\Omega}(X_{ij}-M_{ij})^2}_{\text{只在已知位置计算误差}}
 +\lambda\underbrace{\norm{X}_*}_{\text{促进低秩}}.
\]
误差项使补全接近已知评分，正则项偏好低秩结构。缺失位置没有观测，不能把它们当作零分加入误差。

低秩还可以通过因子表达：$X=UV^{\mathsf T}$，其中 $U\in\R^{m\times r}$、$V\in\R^{n\times r}$。$U$ 的一行描述一个用户在 $r$ 个潜在因素上的倾向，$V$ 的一行描述一部电影的相应属性，两行的内积给出预测评分。由此 $\rank X\le r$。对大规模矩阵，若 $r\ll\min(m,n)$，才可把参数量从 $mn$ 显著降为约 $(m+n)r$；前面的 $2\times2$ 小例用于解释补全，并不体现这种参数节省。

可对因子直接拟合已知评分，并惩罚因子过大：
\[
 \min_{U,V}\sum_{(i,j)\in\Omega}\bigl((UV^{\mathsf T})_{ij}-M_{ij}\bigr)^2
 +\alpha\norm{U}_F^2+\beta\norm{V}_F^2.
\]
$\norm{U}_F^2$ 是全部元素的平方和，$\alpha,\beta>0$。存储变少，但 $U,V$ 相乘，联合问题通常非凸，因子表示也不唯一；正则项并不消除所有不唯一性。

\paragraph{要点归纳}
稀疏恢复从多种符合测量的信号中，依据“只有少量非零分量”选解；低秩补全从多种符合已知评分的矩阵中，依据“少数共同因素能够解释偏好”选解。这些在求解前提出的结构假设，称为先验。$\ell_1$ 与核范数分别作为较易处理的替代目标，促进稀疏与低秩，但正确恢复都需要额外条件。

\section{逻辑回归与支持向量机}\label{sec:classification}
\subsection{逻辑回归：先预测概率}
分类需要输出猫或狗等类别，而线性打分本身可以取任意实数。逻辑回归先对样本计算一个评分，再把评分转换成类别概率。设类别标签为 $b\in\{-1,1\}$，样本特征为 $\va$，则
\[
 s(\va)=\vw^{\mathsf T}\va+c,\qquad
 p=\frac1{1+e^{-s(\va)}}.
\]
权重 $\vw$ 与截距 $c$ 待训练，$s$ 是模型评分，$p$ 是模型给出的正类概率；负类概率为 $1-p$。评分越大，$p$ 越接近 $1$；评分为零时，$p=1/2$。因此以 $1/2$ 为概率阈值，对应的分界面就是 $\vw^{\mathsf T}\va+c=0$。

\begin{example}[从评分读到损失]
本例以及本节损失公式中的 $\log$ 均为自然对数。
若某样本评分为 $s=\log3$，则 $p=1/(1+1/3)=3/4$。若其真实标签是正类，模型给真实类别的概率为 $3/4$，该样本负对数损失为 $-\log(3/4)$，约 $0.288$；若同样评分对应的真实标签是负类，真实类别的概率只有 $1/4$，损失为 $-\log(1/4)$，约 $1.386$。\textbf{给真实类别的概率越低，惩罚越大。}
\end{example}

训练时对全部样本计算这种损失并求和。独立样本下，最小化负对数似然就是
\begin{equation}\label{eq:logistic}
 \boxed{\min_{\vw,c}\sum_{i=1}^N\log\bigl(1+\exp(-b_i s(\va_i))\bigr).}
\end{equation}
其中 $b_i s(\va_i)$ 把正负标签统一进一个式子：分对且评分越有利，损失越小。还可加入平方 $\ell_2$ 或 $\ell_1$ 正则项，分别控制权重或选择特征。

若用 $y_i\in\{0,1\}$ 表示标签，同一损失写成
\[
 -\sum_i\bigl[y_i\log p_i+(1-y_i)\log(1-p_i)\bigr].
\]
$y_i=1$ 时只留下 $-\log p_i$，$y_i=0$ 时只留下 $-\log(1-p_i)$，所以这个式子仍是在惩罚“真实类别的概率太低”。它称为二分类交叉熵损失，这里用自然对数。

虽然名称中有“回归”，逻辑回归通常用于分类；它与预测连续数值的线性回归不同。无正则且训练数据可完全线性分离时，参数还可能不断增大以降低损失，只有下确界而没有有限参数最优解。因此训练问题是否取得最小值，也需要条件。

\subsection{支持向量机：给分界留出余量}
能够分开训练点的边界可能有很多。假设一条数轴上，负类样本在 $0,1$，正类样本在 $3,4$。把分界放在 $1.1$ 或 $2$ 都能分对，但前者离最近样本只有 $0.1$，后者离最近样本为 $1$。后者为输入的小变化留下了更大余量。\textbf{支持向量机（SVM）选择最小间隔尽量大的边界。}

多维情形的分界面写为 $\vw^{\mathsf T}\va+c=0$，$\vw\ne0$，点到分界面的距离为
\[
 \frac{|\vw^{\mathsf T}\va_i+c|}{\norm{\vw}_2}.
\]
分子是评分的绝对值，除以权重长度才得到几何距离。带标签评分 $b_i(\vw^{\mathsf T}\va_i+c)>0$ 表示分在正确一侧。对 $\vw,c$ 同时乘正数不改变分界面，因此在线性可分时可统一规定最小带标签评分为 $1$。

在这一尺度下，所有样本满足 $b_i(\vw^{\mathsf T}\va_i+c)\ge1$，最近样本到分界面的距离是 $1/\norm{\vw}_2$，两侧支撑平面总宽度是 $2/\norm{\vw}_2$。要让间隔大，就要让权重长度小，由此得到\textbf{硬间隔 SVM}：
\[
 \boxed{\begin{aligned}
 \min_{\vw,c}\quad&\frac12\norm{\vw}_2^2,\\
 \text{s.t.}\quad&b_i(\vw^{\mathsf T}\va_i+c)\ge1,\quad i=1,\ldots,N.
 \end{aligned}}
\]
约束要求全部样本分类正确且留出规范后的间隔，目标则把几何间隔尽量扩大。硬间隔中的支持向量是位于间隔边界的样本，限制着边界还能向哪里移动。数轴例中，靠内的 $1$ 与 $3$ 起这种作用。间隔有助于提高一定程度扰动下的稳健性，并不保证任意新数据都分类正确。

\subsection{软间隔：允许违反，但计入代价}
数据重叠或含异常点时，要求每个点都满足间隔可能无法做到。\textbf{软间隔}允许样本欠缺间隔，但要为欠缺量付出代价。令
\[
 m_i=b_i(\vw^{\mathsf T}\va_i+c).
\]
$m_i$ 是带标签评分，\textbf{不是已经除以权重长度的几何距离}。规范要求为 $m_i\ge1$；若允许欠缺 $\xi_i$，就改成 $m_i\ge1-\xi_i$，并要求 $\xi_i\ge0$。

\[
 \boxed{\begin{aligned}
 \min_{\vw,c,\boldsymbol\xi}\quad&\frac12\norm{\vw}_2^2+C\sum_i\xi_i,\\
 \text{s.t.}\quad&m_i\ge1-\xi_i,\qquad\xi_i\ge0,\qquad C>0.
 \end{aligned}}
\]
第一项偏好较大间隔，第二项惩罚间隔不足。对固定的 $\vw,c$，$C>0$ 使算法选择尽量小的可行 $\xi_i$，因此最优欠缺量为
\[
 \xi_i^*=\max\{1-m_i,0\}.
\]
读成中文就是“$1$ 减去带标签评分，若为负数则按零算”。
\begin{center}
\begin{tabular}{ccl}
\toprule
带标签评分 $m_i$ & 最小欠缺量 $\xi_i^*$ & 样本状态\\
\midrule
$1.2$ & $0$ & 分类正确，满足间隔。\\
$0.5$ & $0.5$ & 分类正确，但离分界还不够远。\\
$0$ & $1$ & 位于分界面。\\
$-0.2$ & $1.2$ & 在错误一侧，被误分类。\\
\bottomrule
\end{tabular}
\end{center}
所以松弛量大于零不一定误分类，松弛量之和也不是误分类个数。任意可行 $\xi_i$ 可以取得更大，但解释最小违反程度时应使用上述最优值。

欠缺量函数 $\max\{1-m_i,0\}$ 称为\textbf{合页损失}。消去辅助变量，软间隔模型等价于
\[
 \min_{\vw,c}\frac12\norm{\vw}_2^2
 +C\sum_i\max\{1-b_i(\vw^{\mathsf T}\va_i+c),0\}.
\]
$C$ 越大，违反间隔的代价相对越高，越重视符合训练间隔；$C$ 较小则更容许违反，以换取较小权重。选择 $C$ 也需要验证。若把第一项改为 $\norm{\vw}_1$，得到稀疏权重的变形
\[
 \min_{\vw,c}\norm{\vw}_1+C\sum_i\max\{1-m_i,0\},
\]
它促进特征选择，但已经改变了欧氏平方范数的最大间隔模型。

\paragraph{要点归纳}
逻辑回归从类别概率出发，通过最大似然学习参数；SVM 从分界余量出发，通过最大间隔和违反代价选择边界。逻辑回归输出模型概率，SVM 原始评分不直接等于概率。软间隔同时权衡间隔和训练违反，不能把它简化成“只数错了几个”。两种方法都可以正则化，应按数据与验证表现选择。

\section{优化算法的基本思想}\label{sec:algorithms}
\subsection{方向、步长与一次更新}
目标函数告诉我们什么结果更好，算法则规定怎样从一个起点逐步寻找更好的结果。一般迭代形式为
\[
 \vx^{k+1}=\vx^k+\alpha_k\boldsymbol d^k,
\]
上标 $k$ 表示第几次迭代，$\boldsymbol d^k$ 是搜索方向，$\alpha_k>0$ 是步长。方向表达往哪里走，步长决定走多远。步长太大可能越过有利区域甚至发散，太小则可能进展缓慢。

算法应与模型结构相适应。可微目标能够使用梯度信息，含 $\ell_1$ 范数或合页损失的目标则不能在所有点直接套用普通梯度；整数约束也不能仅靠连续的小幅移动处理。即使数学模型正确，计算过程和停止判断仍需分别检验。

对一元可微函数，导数表示当前位置的变化趋势；多元情况下，各参数的偏导数组成\textbf{梯度}：
\[
 \nabla f(\vx)=\left(\frac{\partial f}{\partial x_1},\ldots,
 \frac{\partial f}{\partial x_n}\right)^{\mathsf T}.
\]
当梯度非零时，沿负梯度方向取充分小的步长可使目标下降。\textbf{梯度下降}采用
\[
 \boxed{\vx^{k+1}=\vx^k-\alpha_k\nabla f(\vx^k).}
\]
它利用一阶变化信息改善目标，并不是直接知道了最优点的位置。

\begin{example}[一次梯度更新]
设目标 $f(x)=(x-3)^2$，从 $x^0=0$ 开始。导数 $f'(x)=2(x-3)$，所以当前位置的导数是 $-6$，负导数方向向右。若步长为 $0.1$，则
\[
 x^1=0-0.1\times(-6)=0.6.
\]
目标由 $f(0)=9$ 降为 $f(0.6)=5.76$，向最小点 $3$ 前进了一步。步长并非越大越好：取 $1$ 会从 $0$ 跳到 $6$，再跳回 $0$；取大于 $1$ 的固定步长，在此例中还会使误差越来越大。
\end{example}

\subsection{全批量、随机与小批量}

大规模学习的目标常为 $f(\vx)=N^{-1}\sum_{i=1}^N f_i(\vx)$。全梯度每次使用全部样本，计算较准确但代价较大。\textbf{随机梯度下降（SGD）}每次随机抽取样本，用单个或少量样本的梯度近似全梯度，例如
\[
 \boldsymbol g^k=\frac1{|B_k|}\sum_{i\in B_k}\nabla f_i(\vx^k),\qquad
 \vx^{k+1}=\vx^k-\alpha_k\boldsymbol g^k.
\]
这里 $B_k$ 是本次抽到的样本集合，$|B_k|$ 是其大小。$\boldsymbol g^k$ 只计算这些样本的平均梯度，然后用它代替全梯度。取一个样本就是单样本随机更新，取若干样本就是小批量更新，取全部样本则回到全批量。适当的均匀抽样下，这种估计平均而言等于全梯度，即无偏；但单次方向可能有误差，完整目标未必每步都下降。

随机方法以带噪声的方向换取较低的单步代价，适合样本量很大的问题。较大的批量通常减少梯度噪声，却增加单步计算量；总体效率还受并行计算、步长与目标结构影响。

在神经网络中，\textbf{前向传播}计算预测和损失，\textbf{反向传播}利用链式法则高效计算各层参数的梯度，优化器再用梯度更新参数。反向传播负责求导，SGD、Adam 等负责更新，不能把二者看作同一个环节。

\subsection{动量、自适应步长与其他方法}
大规模训练中，常在随机梯度基础上积累方向信息或调整各分量的步长。\textbf{动量方法}积累近期方向，减少更新的反复摆动；Nesterov 类方法进一步利用预测位置的信息。\textbf{自适应方法}根据历史梯度大小，对不同坐标采用不同尺度的更新。
\begin{center}
\begin{tabular}{p{0.36\textwidth}p{0.53\textwidth}}
\toprule
方法 & 基本认识\\
\midrule
SGD、ASGD & 随机梯度更新；ASGD 还对迭代结果作平均。\\
Momentum、Nesterov & 利用方向积累或预测位置改善更新。\\
Adagrad & 累积各坐标梯度平方，调整相应学习率。\\
RMSprop、Adadelta & 利用近期平方梯度等尺度信息，避免只按全部历史累积。\\
Adam、Adamax & 结合方向平均与尺度调整；Adamax 使用不同的尺度估计。\\
\bottomrule
\end{tabular}
\end{center}
这些方法改变的是训练目标的求解过程，不会自动修正数据偏差，也不普遍保证找到非凸目标的全局最优解。其他算法名称如 FTRL、ProximalAdagrad、YellowFin，对应不同的结构处理或参数调整思路；理解它们属于优化器及其变种即可，不应把名称列表当作算法推导。

\textbf{牛顿法}还利用坡度怎样变化的信息。一阶近似描述局部的倾斜，二阶近似进一步描述弯曲；牛顿法用局部二次模型来选择方向。对二阶可微的无约束目标，其方向满足
\[
 \nabla^2 f(\vx^k)\boldsymbol d^k=-\nabla f(\vx^k),
\]
其中 $\nabla^2f$ 称为 Hessian 矩阵，收集各参数之间的二阶变化关系。相比只用梯度，计算和求解这组方程通常更贵；在适当条件且靠近合适的极小点时，牛顿法可以很快，但非凸问题中的原始方向未必下降，常需修正或控制步长。

\textbf{原始--对偶方法}同时处理原决策变量与约束对应的乘子。原变量是投资比例等真正要求的方案，乘子是反映各约束作用的辅助变量，也称对偶变量。算法一方面更新方案，另一方面结合约束满足情况更新乘子，二者相互影响，协调目标与约束；它不是简单地预先固定一个罚价后永远不变。具体更新与收敛条件依问题而定。

\paragraph{要点归纳}
梯度法利用局部一阶信息，随机和小批量方法减少每次所用样本，动量积累方向，自适应方法调整不同参数的步长，牛顿法利用二阶关系。比较算法时，既要看一次更新能改善多少，也要看这一步用了多少计算和存储；算法名称本身不能保证预测准确或全局最优。

\section{最优性、收敛与停止准则}\label{sec:convergence}
\subsection{全局最优与局部最优}
局部最优只比较邻近可行点，全局最优则比较全部可行点。可以把目标值想成高度：走进一个小山谷，附近没有更低处，只说明局部较低；若远处还有更深的山谷，当前位置仍不是全局最低。比较时始终只在满足约束的方案中进行。
\begin{definition}
设 $\overline{\vx}\in\mathcal X$。若存在 $\delta>0$，使对所有满足 $\norm{\vx-\overline{\vx}}_2<\delta$ 的可行点都有
\[
 f(\overline{\vx})\le f(\vx),
\]
则称 $\overline{\vx}$ 为局部极小解。若对其中一切 $\vx\ne\overline{\vx}$ 都为严格不等式，则称为严格局部极小解。
\end{definition}
全局极小解一定是局部极小解，反之一般不成立。对凸优化问题，若存在更好的可行点，沿它与局部极小点之间的线段就能在任意近处找到更好的点，与局部极小矛盾，故局部极小必为全局极小。

在可微无约束问题中，梯度为零的点称为\textbf{驻点}。内部极小点一定是驻点，但驻点不一定是极小点，也可能是极大点或鞍点。例如 $f(x)=x^3$ 在 $0$ 处导数为零，却不是极小点。对带约束问题，最优点的普通梯度未必为零，例如 $\min_{x\ge0}x$ 的最优点为 $0$，梯度始终为 $1$。

\subsection{收敛对象与初始点}
评价迭代算法，既要说明什么量趋于什么，也要说明对初始点有什么要求。
\begin{itemize}
\item \textbf{点列收敛：}$\norm{\vx^k-\vx^*}\to0$，迭代点趋向某个解。
\item \textbf{函数值收敛：}$f(\vx^k)\to f^*$，目标值趋向所要求的极小值。
\item \textbf{最优性度量收敛：}例如可微无约束问题中 $\norm{\nabla f(\vx^k)}\to0$，或约束问题中适当的最优性条件违反度趋于零。
\end{itemize}
这三种结论回答不同的问题：参数是否稳定到一个点，目标是否接近所需数值，检查最优性的指标是否变小。若目标连续，点列趋于 $\vx^*$ 可推出函数值趋于 $f(\vx^*)$；反之一般不成立。函数在多个点有同一个最优值时，点列可以在这些点之间移动。梯度范数趋于零也只说明一阶条件的残差趋于零，不直接保证离真实驻点很近。

聚点是某个子列的极限。若梯度范数趋于零，点列存在聚点且梯度连续，则该聚点为驻点，但仍不保证它是极小点或全局最优解。“每个聚点都满足某种最优性条件”也弱于“整个点列收敛到一个最优解”。若迭代过程允许不可行点，还须检查约束违反度能否随迭代减小。

算法的\textbf{局部收敛}通常要求初始点足够接近目标解，\textbf{全局收敛}则在规定的初始点范围内不需要这种邻近要求。这里的“全局”描述初始点条件；\textbf{收敛到全局最优解}描述极限解的性质，两者不能混淆。一个非凸算法即使从任意允许初始点都趋向驻点，也不能因此称其必能求得全局最优解。

\subsection{收敛速度}
在迭代确实收敛的前提下，还可以比较误差衰减的速度。设 $e_k=\norm{\vx^k-\vx^*}>0$，它表示第 $k$ 步离极限点有多远。若误差每步减半，$0.1,0.05,0.025,\ldots$ 是固定比例缩小；若误差每步平方，$0.1,0.01,0.0001,\ldots$ 在已经较小时会缩得更快。正式定义考察 $e_k\to0$ 后充分大的 $k$：
\begin{itemize}
\item \textbf{Q-线性收敛：}存在 $a\in(0,1)$，使 $e_{k+1}\le a e_k$。误差每一步至少按固定比例缩小。
\item \textbf{Q-超线性收敛：}$e_{k+1}/e_k\to0$，误差缩小比例趋于零。
\item \textbf{Q-二次收敛：}存在 $M>0$，使 $e_{k+1}\le M e_k^2$，即下一步误差不超过当前误差平方的某个固定倍数。
\item \textbf{Q-次线性收敛：}在这里按课件的比值表述，$e_{k+1}/e_k\to1$，没有固定的小于 $1$ 的渐近收缩比例。
\end{itemize}
例如 $e_k=2^{-k}$ 为线性收敛，$e_k=2^{-2^k}$ 为二次收敛且超线性收敛，$e_k=1/k$ 为次线性收敛。这里“线性”描述误差缩小规律，不表示误差等于 $ak+b$。

若存在 Q-线性趋于零的非负序列 $t_k$，使 $e_k\le t_k$，则称为\textbf{R-线性收敛}。它用一个收敛较快的上界控制误差，允许原误差序列出现一定波动；R-超线性和 R-二次也可类似理解。

\textbf{迭代复杂度}则研究为保证达到给定精度，需要多少次迭代。例如已有保证
\[
 f(\vx^k)-f^*\le\frac{C_0}{\sqrt{k}},\qquad C_0>0,
\]
为保证误差不超过 $\varepsilon$，取 $k\ge C_0^2/\varepsilon^2$ 即可，因而得到 $O(\varepsilon^{-2})$ 的迭代次数上界。这里 $O$ 表示上界的增长量级：精度要求变为原来的十分之一，这一次数保证约按一百倍增长。这是充分保证，并非每次实际计算都至少需要这么多步。总运行时间还取决于单步计算、内存和通信成本，不能只比较收敛阶。

\subsection{有限步停止}
实际计算必须在有限步结束，因此需要用可计算的量判断是否已经达到所需精度。最优点通常事先未知，不能直接说“离答案很近就停止”。若最优值 $f^*$ 已知，可参考归一化的目标差
\[
 \frac{|f(\vx^k)-f^*|}{\max\{1,|f^*|\}}\le\varepsilon.
\]
但实际中 $f^*$ 往往未知，因此不能把它作为通用的可执行停止条件。对光滑无约束问题，梯度范数较小表示近似满足一阶驻点条件，并不直接表示到真实驻点的距离较小；对约束问题，还需检验可行性，例如
\[
 \max\{g_i(\vx^k),0\}\le\varepsilon_g,\qquad
 |h_j(\vx^k)|\le\varepsilon_h,
\]
并结合适合该问题的最优性条件，不能只检查普通梯度。

常见的停机指标还有
\[
 \frac{\norm{\vx^{k+1}-\vx^k}}{\max\{1,\norm{\vx^k}\}}\le\varepsilon_x,
 \qquad
 \frac{|f(\vx^{k+1})-f(\vx^k)|}{\max\{1,|f(\vx^k)|\}}\le\varepsilon_f.
\]
它们说明迭代或目标变化已经很小，也可能只是步长过小或算法停滞。例如前面的 $f(x)=(x-3)^2$，从 $0$ 出发若只取步长 $10^{-9}$，第一步仅移动 $6\times10^{-9}$；位置几乎没变，但仍远离最小点 $3$，目标值仍接近 $9$。所以不能仅凭变化很小就断言已经最优。
这些比值的分母用于按当前数值的尺度比较变化；与 $1$ 取较大者，也避免当前值为零时无法相除。

\paragraph{要点归纳}
\textbf{停止计算、满足近似最优性条件和找到全局最优解，是不同层次的结论。}梯度和约束检验各有适用条件，资源限制只规定何时结束计算；理论收敛则要说明初始点范围、趋近的对象及所需假设。只有这些关系明确，才能评价一个算法实际得到了什么结果。

\section{主要概念与模型的联系}
\subsection{从数据到决策的知识结构}
数据、学习模型和优化算法分别承担不同任务，又通过同一个问题相互连接。
\begin{center}
\begin{tabular}{p{0.18\textwidth}p{0.34\textwidth}p{0.34\textwidth}}
\toprule
层次 & 核心作用 & 与下一层的联系\\
\midrule
数据与信息 & 记录现象，减少有关目标的不确定性 & 表示与质量决定模型能利用什么。\\
统计与学习 & 从样本中估计规律，预测未知结果 & 通过损失和先验把估计写为优化。\\
优化建模 & 确定变量、目标、约束与结构 & 模型性质决定可采用的算法。\\
迭代求解 & 用有限计算得到近似解 & 检验可行性、最优性和计算代价。\\
验证与决策 & 比较新数据上的效果，选择行动 & 实际反馈用于改进数据与模型。\\
\bottomrule
\end{tabular}
\end{center}
预测能力、数学求解与实际决策效果需要分别评价。例如，一个推荐模型即使准确求解了训练目标，也可能由于样本偏差而推荐不佳；一个训练误差略大的正则化模型，反而可能在新用户上更稳定。

\subsection{模型之间的比较}
比较模型时，需要同时说明输出、损失、结构假设与求解性质。
\begin{center}
\small
\begin{tabular}{p{0.18\textwidth}p{0.30\textwidth}p{0.38\textwidth}}
\toprule
模型 & 主要任务与损失 & 结构要求或主要区别\\
\midrule
最小二乘 & 连续响应；残差平方和 & 控制整体平方误差，不自动进行特征选择。\\
岭回归 & 平方损失加 $\ell_2$ 平方 & 缩小系数、改善稳定性，一般不产生大量精确零。\\
LASSO & 平方损失加 $\ell_1$ & 促进系数稀疏，涉及非光滑目标。\\
逻辑回归 & 二分类；负对数似然 & 用概率解释输出，线性评分下边界为超平面。\\
软间隔 SVM & 二分类；合页损失 & 权衡权重范数与间隔违反程度。\\
核范数补全 & 预测矩阵缺失元素 & 以奇异值之和促进低秩，仅在观测位置拟合。\\
低秩因子模型 & 矩阵拟合与因子正则 & 大矩阵、小秩时节省参数，联合优化通常非凸。\\
\bottomrule
\end{tabular}
\end{center}

若干相近概念应保持各自含义：数据量不等于有效信息量，相关关系不等于因果关系；观测稀疏不等于矩阵低秩，参数缩小不等于特征筛选；模型是比较方案的数学表达，算法是计算方案；优化目标收敛不等于新数据预测正确，停机不等于最优。

完整说明一种方法，应使问题、原理、表达、案例与适用条件相互对应。对于大数据应用，说明记录如何形成信息、信息如何支持判断；对于优化模型，说明每个变量、目标项和约束的实际意义；对于算法，说明所用信息、更新思想与计算特点，再区分可证明的保证和实际停止条件。这样才能把概念、公式和案例连成完整解释。

从记录中取得信息，用模型形成预测，再结合目标与约束选择行动，最后通过新的观测检验结果，数据分析与最优化就在这一过程中相互配合。

\smallskip
\noindent\textbf{答题安排：}6 道论述题共 100 分钟，可以先用约 5 分钟通读题目、按分值分配时间，保留约 5 分钟检查。若各题分值相近，正文作答可按每题约 15 分钟安排；具体随分值和内容调整。作答以切合题意的定义与关系为起点，展开机制，选用相应案例或公式，最后交代成立条件，避免只罗列术语。

\clearpage
\appendix
\section{课件内容索引}
以下均为 PDF 页码。正文按知识关系重组，重复内容合并，历史数据与软件链接不作为原理的记忆要求。
\begin{longtable}{p{0.25\textwidth}p{0.15\textwidth}p{0.47\textwidth}}
\toprule
课件 & 页码 & 对应内容\\
\midrule
\endhead
大数据案例介绍 & 8--24 & 数据、信息、知识与 4V；数据与模型。\\
 & 25--39 & 推荐、流感、药物、广告与行为记录。\\
 & 40--44 & 样本量、代表性与数据驱动。\\
 & 45--54 & 博弈、运筹学及与数据科学的联系。\\
 & 55--60 & 数据价值、行为推荐、社会网络、属性预测与隐私。\\
\midrule
大数据思维 & 3--13 & 大数据特点、多维性、翻译与交通案例。\\
 & 14--18 & 确定性、不确定性、信息熵、编码、最大熵与交叉熵。\\
 & 19--26 & 异常检测、机器智能、语音与问答。\\
 & 27--33 & 知识库、机器学习、表示学习与深度学习。\\
\midrule
最优化：简介 & 3--9 & 优化的一般形式、分类与投资组合。\\
 & 11--17 & 稀疏优化、基追踪与 LASSO。\\
 & 19--23 & 低秩恢复、核范数与因子分解。\\
 & 25--35 & 学习模型、MLP、CNN、转置卷积、RNN 与优化器。\\
 & 37--45 & 最优解、收敛、速度、复杂度与停止准则。\\
\midrule
优化建模 & 2--6 & 建模流程、误差度量与正则化。\\
 & 7--10 & 最大似然、代价收益及替代目标。\\
 & 11--14 & 约束、等价变换与松弛。\\
 & 15--24 & 一般学习目标、回归、岭、LASSO 及结构变形。\\
 & 25--27 & 逻辑回归、期望风险与样本平均。\\
 & 28--31 & 硬间隔、软间隔支持向量机及合页损失。\\
\bottomrule
\end{longtable}
\end{document}
